Thus @@M0@@, which makes @@M1@@ an @@M2@@-torsor over @@M3@@, is affine. Moreover, it is known (SGA 3 1.4.2) that @@M4@@ is a closed immersion, hence affine, and hence the same holds for the composite @@M5@@.

\sgaStableTarget{sga6.e12.cor.1-5}\textbf{Corollary 1.5.} \uline{Let} @@M6@@ \uline{be a field.}
\begin{enumerate}[label=\alph*)]
\item \uline{Let} @@M7@@ \uline{be a proper} @@M8@@\uline{-scheme. Then} @@M9@@ \uline{is representable} (J. P. Murre).
\item \uline{Let} @@M10@@ \uline{be a surjective morphism between two proper} @@M11@@\uline{-schemes. Then the induced morphism} @@M12@@ \uline{is affine.}
\end{enumerate}

\sgaStableTarget{sga6.e12.cor.1-6}\textbf{Corollary 1.6.} \uline{Let} @@M13@@ \uline{be an integral noetherian scheme, and} @@M14@@ \uline{a surjective} @@M15@@\uline{-morphism between two proper} @@M16@@\uline{-schemes.}
\begin{enumerate}[label=\alph*)]
\item \uline{If} @@M17@@ \uline{is normal over} @@M18@@ \uline{(that is, flat with geometrically normal fibers), then, after shrinking} @@M19@@, @@M20@@ \uline{is quasi-finite.}
\item \uline{If} @@M21@@ \uline{is essentially proper over} @@M22@@ \uline{(that is, if it satisfies the valuative criterion of properness, for example, if} @@M23@@ \uline{is smooth over} @@M24@@\uline{), then, after shrinking} @@M25@@, @@M26@@ \uline{is finite (and in particular affine).}
\end{enumerate}

Indeed, in a), for every @@M27@@, @@M28@@ is an affine algebraic group; its identity component is a proper algebraic group over @@M29@@ (FGA 236, Theorem 2.1 (iii)); @@M30@@ is therefore finite. In b), @@M31@@ is both proper and quasi-affine, hence finite.

\sgaStableTarget{sga6.e12.rem.1-7}\textbf{Remark 1.7.} It is not known whether 1.1 remains valid when ``quasi-affine'' is replaced there by affine, as 1.5 and 1.6 might suggest.

